% algo_flow.tex — Algorithm flowchart: Klempt 2024 Eq.34-36 + UMAT coupling
% v2: fixed text overlaps (label positions, loop arrow note, coupling note)
%
% Usage (thesis or Beamer):
%   \resizebox{\linewidth}{!}{% algo_flow.tex — Algorithm flowchart: Klempt 2024 Eq.34-36 + UMAT coupling
% v2: fixed text overlaps (label positions, loop arrow note, coupling note)
%
% Usage (thesis or Beamer):
%   \resizebox{\linewidth}{!}{% algo_flow.tex — Algorithm flowchart: Klempt 2024 Eq.34-36 + UMAT coupling
% v2: fixed text overlaps (label positions, loop arrow note, coupling note)
%
% Usage (thesis or Beamer):
%   \resizebox{\linewidth}{!}{% algo_flow.tex — Algorithm flowchart: Klempt 2024 Eq.34-36 + UMAT coupling
% v2: fixed text overlaps (label positions, loop arrow note, coupling note)
%
% Usage (thesis or Beamer):
%   \resizebox{\linewidth}{!}{\input{algo_flow.tex}}
% ─────────────────────────────────────────────────────────────────────
\begin{tikzpicture}[
  every node/.style = {font=\small},
  box/.style   = {draw, rounded corners=4pt, minimum width=7.2cm,
                  minimum height=0.9cm, align=center, thick},
  init/.style  = {box, fill=gray!15},
  pde/.style   = {box, fill=blue!12},
  mech/.style  = {box, fill=orange!18},
  pass/.style  = {box, fill=green!18},
  arrow/.style = {-Stealth, thick},
  loopnote/.style = {font=\scriptsize\itshape, text width=2.6cm, align=right},
  cnote/.style    = {font=\scriptsize\itshape, text width=3.0cm, align=center,
                     fill=white, inner sep=1.5pt},
]

%% ── Column A: PDE time loop ──────────────────────────────────────────
\node[init] (init)
  {Initialise \;$\phi^0$,\; $c^0$,\; $\alpha^0 = 0$};

\node[pde, below=0.55cm of init] (pois)
  {Solve nutrient (quasi-static), Eq.\,35:\\
   $-D\,\Delta c + g\,\phi = 0$};

\node[pde, below=0.55cm of pois] (phi)
  {Update $\phi$ (Eq.\,34, explicit Euler):\\[1pt]
   $\phi^{n+1} = \phi^n + \Delta t\!\left[\,
     \beta\Delta\phi
     - k_\alpha\alpha
     + \frac{rc}{k{+}c}\,(\nabla\phi\cdot\hat n_c)
   \right]$};

\node[pde, below=0.55cm of phi] (alp)
  {Update growth variable (Eq.\,36):\\[1pt]
   $\alpha^{n+1} = \alpha^n + \Delta t\,k_\alpha\,\phi$};

\node[init, below=0.55cm of alp] (loop)
  {$t \leftarrow t + \Delta t$ \quad Repeat until $t = T_{\max}$};

%% ── Column B: UMAT / Gauss-point coupling ────────────────────────────
\node[mech, right=2.8cm of pois, yshift=-1.6cm] (growth)
  {Growth split \;$\mathbf{F} = \mathbf{F}_e\cdot\mathbf{F}_g$\\[1pt]
   $\mathbf{F}_g = \mathrm{diag}(1,1,\lambda_z)$,\;
   $\mathbf{F}_e = \mathbf{F}\,\mathbf{F}_g^{-1}$\\[-1pt]
   {\footnotesize$\lambda_z = 1 + \max(0,\,\beta_d\,\phi_{Pg} + b)$}};

\node[mech, below=0.55cm of growth] (stress)
  {neo-Hookean Cauchy stress,\; $J{=}\det\mathbf{F}_e$\\[1pt]
   $\boldsymbol\sigma
     = \dfrac{\mu}{J}(\mathbf{b}_e - \mathbf{I})
     + \dfrac{\lambda}{J}\ln J\;\mathbf{I}$};

\node[mech, below=0.55cm of stress] (tang)
  {Consistent spatial tangent (Bonet \& Wood §6.5)\\[1pt]
   $\mathbb{C}_{ijkl}
     = \dfrac{\lambda}{J}\delta_{ij}\delta_{kl}
     + \dfrac{\mu {-} \lambda\ln J}{J}
       (\delta_{ik}\delta_{jl} + \delta_{il}\delta_{jk})$};

\node[pass, below=0.55cm of tang] (ret)
  {\textbf{Return} $(\boldsymbol\sigma,\,\mathbb{C})$ to Abaqus Newton loop};

%% ── Arrows: PDE loop ─────────────────────────────────────────────────
\draw[arrow] (init) -- (pois);
\draw[arrow] (pois) -- (phi);
\draw[arrow] (phi)  -- (alp);
\draw[arrow] (alp)  -- (loop);
% back-edge: place label on the LEFT stub, not on the bend
\draw[arrow] (loop.west) -- ++(-0.85,0) coordinate (lstub)
             |- (pois.west);
\node[loopnote, left=0.05cm of lstub] {every $C_{\mathrm{solve}}$ steps};

%% ── Arrow: PDE → UMAT coupling ───────────────────────────────────────
% go right from phi, then down-right to growth; label on a white bg node
\draw[arrow, dashed, red!70!black, thick]
  (phi.east) -- ++(0.5,0) coordinate (pmid)
             |- (growth.west);
\node[cnote] at ($(pmid)!0.5!(growth.west -| pmid)$)
  {$\phi_{Pg}(x)$ to each\\ Gauss point};

%% ── Arrows: UMAT internal ────────────────────────────────────────────
\draw[arrow] (growth) -- (stress);
\draw[arrow] (stress) -- (tang);
\draw[arrow] (tang)   -- (ret);

%% ── Background boxes (labels placed BELOW top border, inside the box) ─
\begin{scope}[on background layer]
  \node[draw=blue!40, rounded corners=6pt, fill=blue!4,
        fit=(init)(pois)(phi)(alp)(loop), inner sep=8pt] (pdeBox) {};
  \node[draw=orange!50, rounded corners=6pt, fill=orange!4,
        fit=(growth)(stress)(tang)(ret), inner sep=8pt] (umatBox) {};
\end{scope}
% labels outside the box (above, centred), no overlap with content
\node[font=\footnotesize\bfseries, text=blue!60,
      above=3pt of pdeBox.north] {PDE loop\enspace(\texttt{felix\_complete\_reproduction.py})};
\node[font=\footnotesize\bfseries, text=orange!80!black,
      above=3pt of umatBox.north] {UMAT / Gauss-point\enspace(\texttt{umat\_server\_nsp\_fs.py})};

\end{tikzpicture}
}
% ─────────────────────────────────────────────────────────────────────
\begin{tikzpicture}[
  every node/.style = {font=\small},
  box/.style   = {draw, rounded corners=4pt, minimum width=7.2cm,
                  minimum height=0.9cm, align=center, thick},
  init/.style  = {box, fill=gray!15},
  pde/.style   = {box, fill=blue!12},
  mech/.style  = {box, fill=orange!18},
  pass/.style  = {box, fill=green!18},
  arrow/.style = {-Stealth, thick},
  loopnote/.style = {font=\scriptsize\itshape, text width=2.6cm, align=right},
  cnote/.style    = {font=\scriptsize\itshape, text width=3.0cm, align=center,
                     fill=white, inner sep=1.5pt},
]

%% ── Column A: PDE time loop ──────────────────────────────────────────
\node[init] (init)
  {Initialise \;$\phi^0$,\; $c^0$,\; $\alpha^0 = 0$};

\node[pde, below=0.55cm of init] (pois)
  {Solve nutrient (quasi-static), Eq.\,35:\\
   $-D\,\Delta c + g\,\phi = 0$};

\node[pde, below=0.55cm of pois] (phi)
  {Update $\phi$ (Eq.\,34, explicit Euler):\\[1pt]
   $\phi^{n+1} = \phi^n + \Delta t\!\left[\,
     \beta\Delta\phi
     - k_\alpha\alpha
     + \frac{rc}{k{+}c}\,(\nabla\phi\cdot\hat n_c)
   \right]$};

\node[pde, below=0.55cm of phi] (alp)
  {Update growth variable (Eq.\,36):\\[1pt]
   $\alpha^{n+1} = \alpha^n + \Delta t\,k_\alpha\,\phi$};

\node[init, below=0.55cm of alp] (loop)
  {$t \leftarrow t + \Delta t$ \quad Repeat until $t = T_{\max}$};

%% ── Column B: UMAT / Gauss-point coupling ────────────────────────────
\node[mech, right=2.8cm of pois, yshift=-1.6cm] (growth)
  {Growth split \;$\mathbf{F} = \mathbf{F}_e\cdot\mathbf{F}_g$\\[1pt]
   $\mathbf{F}_g = \mathrm{diag}(1,1,\lambda_z)$,\;
   $\mathbf{F}_e = \mathbf{F}\,\mathbf{F}_g^{-1}$\\[-1pt]
   {\footnotesize$\lambda_z = 1 + \max(0,\,\beta_d\,\phi_{Pg} + b)$}};

\node[mech, below=0.55cm of growth] (stress)
  {neo-Hookean Cauchy stress,\; $J{=}\det\mathbf{F}_e$\\[1pt]
   $\boldsymbol\sigma
     = \dfrac{\mu}{J}(\mathbf{b}_e - \mathbf{I})
     + \dfrac{\lambda}{J}\ln J\;\mathbf{I}$};

\node[mech, below=0.55cm of stress] (tang)
  {Consistent spatial tangent (Bonet \& Wood §6.5)\\[1pt]
   $\mathbb{C}_{ijkl}
     = \dfrac{\lambda}{J}\delta_{ij}\delta_{kl}
     + \dfrac{\mu {-} \lambda\ln J}{J}
       (\delta_{ik}\delta_{jl} + \delta_{il}\delta_{jk})$};

\node[pass, below=0.55cm of tang] (ret)
  {\textbf{Return} $(\boldsymbol\sigma,\,\mathbb{C})$ to Abaqus Newton loop};

%% ── Arrows: PDE loop ─────────────────────────────────────────────────
\draw[arrow] (init) -- (pois);
\draw[arrow] (pois) -- (phi);
\draw[arrow] (phi)  -- (alp);
\draw[arrow] (alp)  -- (loop);
% back-edge: place label on the LEFT stub, not on the bend
\draw[arrow] (loop.west) -- ++(-0.85,0) coordinate (lstub)
             |- (pois.west);
\node[loopnote, left=0.05cm of lstub] {every $C_{\mathrm{solve}}$ steps};

%% ── Arrow: PDE → UMAT coupling ───────────────────────────────────────
% go right from phi, then down-right to growth; label on a white bg node
\draw[arrow, dashed, red!70!black, thick]
  (phi.east) -- ++(0.5,0) coordinate (pmid)
             |- (growth.west);
\node[cnote] at ($(pmid)!0.5!(growth.west -| pmid)$)
  {$\phi_{Pg}(x)$ to each\\ Gauss point};

%% ── Arrows: UMAT internal ────────────────────────────────────────────
\draw[arrow] (growth) -- (stress);
\draw[arrow] (stress) -- (tang);
\draw[arrow] (tang)   -- (ret);

%% ── Background boxes (labels placed BELOW top border, inside the box) ─
\begin{scope}[on background layer]
  \node[draw=blue!40, rounded corners=6pt, fill=blue!4,
        fit=(init)(pois)(phi)(alp)(loop), inner sep=8pt] (pdeBox) {};
  \node[draw=orange!50, rounded corners=6pt, fill=orange!4,
        fit=(growth)(stress)(tang)(ret), inner sep=8pt] (umatBox) {};
\end{scope}
% labels outside the box (above, centred), no overlap with content
\node[font=\footnotesize\bfseries, text=blue!60,
      above=3pt of pdeBox.north] {PDE loop\enspace(\texttt{felix\_complete\_reproduction.py})};
\node[font=\footnotesize\bfseries, text=orange!80!black,
      above=3pt of umatBox.north] {UMAT / Gauss-point\enspace(\texttt{umat\_server\_nsp\_fs.py})};

\end{tikzpicture}
}
% ─────────────────────────────────────────────────────────────────────
\begin{tikzpicture}[
  every node/.style = {font=\small},
  box/.style   = {draw, rounded corners=4pt, minimum width=7.2cm,
                  minimum height=0.9cm, align=center, thick},
  init/.style  = {box, fill=gray!15},
  pde/.style   = {box, fill=blue!12},
  mech/.style  = {box, fill=orange!18},
  pass/.style  = {box, fill=green!18},
  arrow/.style = {-Stealth, thick},
  loopnote/.style = {font=\scriptsize\itshape, text width=2.6cm, align=right},
  cnote/.style    = {font=\scriptsize\itshape, text width=3.0cm, align=center,
                     fill=white, inner sep=1.5pt},
]

%% ── Column A: PDE time loop ──────────────────────────────────────────
\node[init] (init)
  {Initialise \;$\phi^0$,\; $c^0$,\; $\alpha^0 = 0$};

\node[pde, below=0.55cm of init] (pois)
  {Solve nutrient (quasi-static), Eq.\,35:\\
   $-D\,\Delta c + g\,\phi = 0$};

\node[pde, below=0.55cm of pois] (phi)
  {Update $\phi$ (Eq.\,34, explicit Euler):\\[1pt]
   $\phi^{n+1} = \phi^n + \Delta t\!\left[\,
     \beta\Delta\phi
     - k_\alpha\alpha
     + \frac{rc}{k{+}c}\,(\nabla\phi\cdot\hat n_c)
   \right]$};

\node[pde, below=0.55cm of phi] (alp)
  {Update growth variable (Eq.\,36):\\[1pt]
   $\alpha^{n+1} = \alpha^n + \Delta t\,k_\alpha\,\phi$};

\node[init, below=0.55cm of alp] (loop)
  {$t \leftarrow t + \Delta t$ \quad Repeat until $t = T_{\max}$};

%% ── Column B: UMAT / Gauss-point coupling ────────────────────────────
\node[mech, right=2.8cm of pois, yshift=-1.6cm] (growth)
  {Growth split \;$\mathbf{F} = \mathbf{F}_e\cdot\mathbf{F}_g$\\[1pt]
   $\mathbf{F}_g = \mathrm{diag}(1,1,\lambda_z)$,\;
   $\mathbf{F}_e = \mathbf{F}\,\mathbf{F}_g^{-1}$\\[-1pt]
   {\footnotesize$\lambda_z = 1 + \max(0,\,\beta_d\,\phi_{Pg} + b)$}};

\node[mech, below=0.55cm of growth] (stress)
  {neo-Hookean Cauchy stress,\; $J{=}\det\mathbf{F}_e$\\[1pt]
   $\boldsymbol\sigma
     = \dfrac{\mu}{J}(\mathbf{b}_e - \mathbf{I})
     + \dfrac{\lambda}{J}\ln J\;\mathbf{I}$};

\node[mech, below=0.55cm of stress] (tang)
  {Consistent spatial tangent (Bonet \& Wood §6.5)\\[1pt]
   $\mathbb{C}_{ijkl}
     = \dfrac{\lambda}{J}\delta_{ij}\delta_{kl}
     + \dfrac{\mu {-} \lambda\ln J}{J}
       (\delta_{ik}\delta_{jl} + \delta_{il}\delta_{jk})$};

\node[pass, below=0.55cm of tang] (ret)
  {\textbf{Return} $(\boldsymbol\sigma,\,\mathbb{C})$ to Abaqus Newton loop};

%% ── Arrows: PDE loop ─────────────────────────────────────────────────
\draw[arrow] (init) -- (pois);
\draw[arrow] (pois) -- (phi);
\draw[arrow] (phi)  -- (alp);
\draw[arrow] (alp)  -- (loop);
% back-edge: place label on the LEFT stub, not on the bend
\draw[arrow] (loop.west) -- ++(-0.85,0) coordinate (lstub)
             |- (pois.west);
\node[loopnote, left=0.05cm of lstub] {every $C_{\mathrm{solve}}$ steps};

%% ── Arrow: PDE → UMAT coupling ───────────────────────────────────────
% go right from phi, then down-right to growth; label on a white bg node
\draw[arrow, dashed, red!70!black, thick]
  (phi.east) -- ++(0.5,0) coordinate (pmid)
             |- (growth.west);
\node[cnote] at ($(pmid)!0.5!(growth.west -| pmid)$)
  {$\phi_{Pg}(x)$ to each\\ Gauss point};

%% ── Arrows: UMAT internal ────────────────────────────────────────────
\draw[arrow] (growth) -- (stress);
\draw[arrow] (stress) -- (tang);
\draw[arrow] (tang)   -- (ret);

%% ── Background boxes (labels placed BELOW top border, inside the box) ─
\begin{scope}[on background layer]
  \node[draw=blue!40, rounded corners=6pt, fill=blue!4,
        fit=(init)(pois)(phi)(alp)(loop), inner sep=8pt] (pdeBox) {};
  \node[draw=orange!50, rounded corners=6pt, fill=orange!4,
        fit=(growth)(stress)(tang)(ret), inner sep=8pt] (umatBox) {};
\end{scope}
% labels outside the box (above, centred), no overlap with content
\node[font=\footnotesize\bfseries, text=blue!60,
      above=3pt of pdeBox.north] {PDE loop\enspace(\texttt{felix\_complete\_reproduction.py})};
\node[font=\footnotesize\bfseries, text=orange!80!black,
      above=3pt of umatBox.north] {UMAT / Gauss-point\enspace(\texttt{umat\_server\_nsp\_fs.py})};

\end{tikzpicture}
}
% ─────────────────────────────────────────────────────────────────────
\begin{tikzpicture}[
  every node/.style = {font=\small},
  box/.style   = {draw, rounded corners=4pt, minimum width=7.2cm,
                  minimum height=0.9cm, align=center, thick},
  init/.style  = {box, fill=gray!15},
  pde/.style   = {box, fill=blue!12},
  mech/.style  = {box, fill=orange!18},
  pass/.style  = {box, fill=green!18},
  arrow/.style = {-Stealth, thick},
  loopnote/.style = {font=\scriptsize\itshape, text width=2.6cm, align=right},
  cnote/.style    = {font=\scriptsize\itshape, text width=3.0cm, align=center,
                     fill=white, inner sep=1.5pt},
]

%% ── Column A: PDE time loop ──────────────────────────────────────────
\node[init] (init)
  {Initialise \;$\phi^0$,\; $c^0$,\; $\alpha^0 = 0$};

\node[pde, below=0.55cm of init] (pois)
  {Solve nutrient (quasi-static), Eq.\,35:\\
   $-D\,\Delta c + g\,\phi = 0$};

\node[pde, below=0.55cm of pois] (phi)
  {Update $\phi$ (Eq.\,34, explicit Euler):\\[1pt]
   $\phi^{n+1} = \phi^n + \Delta t\!\left[\,
     \beta\Delta\phi
     - k_\alpha\alpha
     + \frac{rc}{k{+}c}\,(\nabla\phi\cdot\hat n_c)
   \right]$};

\node[pde, below=0.55cm of phi] (alp)
  {Update growth variable (Eq.\,36):\\[1pt]
   $\alpha^{n+1} = \alpha^n + \Delta t\,k_\alpha\,\phi$};

\node[init, below=0.55cm of alp] (loop)
  {$t \leftarrow t + \Delta t$ \quad Repeat until $t = T_{\max}$};

%% ── Column B: UMAT / Gauss-point coupling ────────────────────────────
\node[mech, right=2.8cm of pois, yshift=-1.6cm] (growth)
  {Growth split \;$\mathbf{F} = \mathbf{F}_e\cdot\mathbf{F}_g$\\[1pt]
   $\mathbf{F}_g = \mathrm{diag}(1,1,\lambda_z)$,\;
   $\mathbf{F}_e = \mathbf{F}\,\mathbf{F}_g^{-1}$\\[-1pt]
   {\footnotesize$\lambda_z = 1 + \max(0,\,\beta_d\,\phi_{Pg} + b)$}};

\node[mech, below=0.55cm of growth] (stress)
  {neo-Hookean Cauchy stress,\; $J{=}\det\mathbf{F}_e$\\[1pt]
   $\boldsymbol\sigma
     = \dfrac{\mu}{J}(\mathbf{b}_e - \mathbf{I})
     + \dfrac{\lambda}{J}\ln J\;\mathbf{I}$};

\node[mech, below=0.55cm of stress] (tang)
  {Consistent spatial tangent (Bonet \& Wood §6.5)\\[1pt]
   $\mathbb{C}_{ijkl}
     = \dfrac{\lambda}{J}\delta_{ij}\delta_{kl}
     + \dfrac{\mu {-} \lambda\ln J}{J}
       (\delta_{ik}\delta_{jl} + \delta_{il}\delta_{jk})$};

\node[pass, below=0.55cm of tang] (ret)
  {\textbf{Return} $(\boldsymbol\sigma,\,\mathbb{C})$ to Abaqus Newton loop};

%% ── Arrows: PDE loop ─────────────────────────────────────────────────
\draw[arrow] (init) -- (pois);
\draw[arrow] (pois) -- (phi);
\draw[arrow] (phi)  -- (alp);
\draw[arrow] (alp)  -- (loop);
% back-edge: place label on the LEFT stub, not on the bend
\draw[arrow] (loop.west) -- ++(-0.85,0) coordinate (lstub)
             |- (pois.west);
\node[loopnote, left=0.05cm of lstub] {every $C_{\mathrm{solve}}$ steps};

%% ── Arrow: PDE → UMAT coupling ───────────────────────────────────────
% go right from phi, then down-right to growth; label on a white bg node
\draw[arrow, dashed, red!70!black, thick]
  (phi.east) -- ++(0.5,0) coordinate (pmid)
             |- (growth.west);
\node[cnote] at ($(pmid)!0.5!(growth.west -| pmid)$)
  {$\phi_{Pg}(x)$ to each\\ Gauss point};

%% ── Arrows: UMAT internal ────────────────────────────────────────────
\draw[arrow] (growth) -- (stress);
\draw[arrow] (stress) -- (tang);
\draw[arrow] (tang)   -- (ret);

%% ── Background boxes (labels placed BELOW top border, inside the box) ─
\begin{scope}[on background layer]
  \node[draw=blue!40, rounded corners=6pt, fill=blue!4,
        fit=(init)(pois)(phi)(alp)(loop), inner sep=8pt] (pdeBox) {};
  \node[draw=orange!50, rounded corners=6pt, fill=orange!4,
        fit=(growth)(stress)(tang)(ret), inner sep=8pt] (umatBox) {};
\end{scope}
% labels outside the box (above, centred), no overlap with content
\node[font=\footnotesize\bfseries, text=blue!60,
      above=3pt of pdeBox.north] {PDE loop\enspace(\texttt{felix\_complete\_reproduction.py})};
\node[font=\footnotesize\bfseries, text=orange!80!black,
      above=3pt of umatBox.north] {UMAT / Gauss-point\enspace(\texttt{umat\_server\_nsp\_fs.py})};

\end{tikzpicture}
